\chapter{Практика}
\label{ch:chap2}

\section*{Задание 1: Для заданных значений частоты синала и частоты дискретизации получите дискретное колебание. Далее
увеличьте частоту сигнала в несколько раз, при этом так же увеличится и частота дискретизации, но отношение частоты сигнала и частоты
дискретизации - нормированная частота останется той же величиной. Сравните дискретные отсчеты первого и второго сигналов.
}

Зададим сигналы с частотой f=32Гц и f=320Гц и частотой дискретизации fs=32*f. В обоих случаях цифровая частота $\Omega = 2\pi\frac{f}{fs} = \frac{1}{32}$,
т.е в цифровом мире эти два сигнала абсолютно идентичны, хотя имеют разные частоты. Выведем в терминал первые 25 отсчетов сигналов
для сравнения

\begin{figure}[H]
    \centering
    \includegraphics[width=1.0\textwidth]{samples.png}
    \caption{Семплы}
\end{figure}

Сверху отсчеты сигнала с f=32Гц, а снизу с f=320Гц. Отсчеты абсолютно идентичны.

\subsection*{Задание 2: Вычислите шаг частот между точками ДПФ $\triangle f =\frac{fs}{N}$. Определите, в какой точке ДПФ находится заданный сигнал}

Сигнал имеет частоту f=32Гц и N=256 точек, тогда $\triangle f =\frac{1024}{256}$ = 4Гц на отсчет. Если наш сигнал имеет частоту
f=32Гц, то он будет иметь номер $\frac{32}{4}$ = 8. Т.е в 8 отсчете ДПФ будет находиться значение компоненты. Для проверки выведем
восьмой элемент массива с ДПФ отсчетами

\begin{figure}[H]
    \centering
    \includegraphics[width=1.0\textwidth]{s8.png}
    \caption{Точка ДПФ}
\end{figure}

Получили число, которое не близко к нулю, значит, это одна из компонент. Поскольку ДПФ симметрично, то такое же значение будет в 
ячейке массива с номером 256-8=248.

\subsection*{Задание 3: Измените частоту сигнала в целое чисто раз, определите номер точки ДПФ для данного сигнала}

Увеличим частоту сигнала в 7 раз, тогда f=224Гц, fs = 1024Гц. $\triangle f = \frac{1024}{256} = 4$. Номер отсчета с компонентой равен
$\frac{224}{4} = 56$. Проверим

\begin{figure}[H]
    \centering
    \includegraphics[width=1.0\textwidth]{s56.png}
    \caption{Точка ДПФ}
\end{figure}

При увеличении частоты в n раз точка ДПФ сместится в отсчетах ДПФ вправо, что логично, ведь частота возросла и теперь алгоритм найдет
эту компоненту позже.

\section*{Задание 4: Измените количество точек ДПФ до 512. Вычислите шаг частот между точками ДПФ. Определите, в какой точке ДПФ находится заданный сигнал}

Если f=32Гц, N=512, то $\triangle f = \frac{1024}{512}$ = 2Гц. Точка ДПФ будет иметь номер $\frac{32}{2}=16$. \\ 
Из этого можно сделать вывод, что при фиксированной частоте дискретизации с увеличением N частотная сетка будет плотнее (шаг будет становиться меньше),
соответственно, нужно будет перебирать больше точек, прежде чем мы попадем в точку ДПФ.


\section*{Задание 5: Задайте сигнал в виде суммы двух колебаний. Вычислите ДПФ сигнала}

Зададим сигнал в виде суммы сигналов с частотами $f_1$=32Гц $f_2$=16Гц. Вычислим ДПФ

\begin{figure}[H]
    \centering
    \includegraphics[width=1.0\textwidth]{spec1.png}
    \caption{Спектр}
\end{figure}

Четко видно 2 спектральные линии на частотах 16Гц и 32Гц.\\

ДПФ принимает только частоты кратные $\triangle f$. В нашем случае подходят частоты кратные 4. Возьмем частоту $f_2=65$

\begin{figure}[H]
    \centering
    \includegraphics[width=1.0\textwidth]{spec2.png}
    \caption{Спектр}
\end{figure}

Получили спектр с четкими спектральными линиями в точках 32 и 64Гц. Энергия от спектарльной линии в точке 65Гц "размазалась" по соседним точкам.

\section*{Задание 6: Вычислите ОДПФ сигнала, заданого в частотной области в виде [0,0,2-1j,0,0,0,0,0], 
поменяйте знак мнимой единицы. Можно увеличить количество точек до 16 при одном ненулевом значении.}

Создадим массив на 16 элементов и заполним его нулями. Элементы 0 и 15 заполним числами 5-10j

\begin{figure}[H]
    \centering
    \includegraphics[width=1.0\textwidth]{sig1.png}
    \caption{Восстановленный сигнал}
\end{figure}

Поменяем числа 5-10j на 5+10j

\begin{figure}[H]
    \centering
    \includegraphics[width=1.0\textwidth]{sig2.png}
    \caption{Восстановленный сигнал}
\end{figure}

Получили тот же сигнал, но отзеркаленый. \\

Сдвинем числа на позиции 3 и 13

\begin{figure}[H]
    \centering
    \includegraphics[width=1.0\textwidth]{sig3.png}
    \caption{Восстановленный сигнал}
\end{figure}

Получили сигнал плохого качества.

\endinput 